\subsection{A form-three class that survives the iterated saturation}

\begin{proposition}[A-S09-22: Proved; the saturated resolved comparison fails in degree three]
\label{AS0922}
Over $R=\mathbb Z_{(2)}$, use the actual resolution and
resolved-tail construction of A-L11, with
$L=I_1\langle s\rangle$, $m=2$, and $s\geq5$.
For the zero twist over $G_R$ there is a class of exact
order two in form degree three and weight $2s+5$, in
simplex two.  It remains nonzero after the full iterated
PD and based $H^1$ saturation of A-SAT01.

The same assertion holds for the actual independent
quadratic class of A-L11: start with the free identity
disk $x_0K(2)\xrightarrow{d}y_0K(2)$ over $G_R$,
with forms one and two and weight $b\geq3$, and put
$s=2b$.  Apply the resolved-tail repair to the coefficient
tail of $z_0=\gamma_2(y_0)$ with
$\eta_0=x_0y_0$, so $d\eta_0=2z_0$.
The class constructed below survives saturation and maps
to zero under the actual saturated comparison to the
two-term tail algebra of A-S09-21.
Both specified target maps may be taken to be the strict
projection to $G_R$.

This is a failure of that particular integral comparison
in degree three.  It does not assert a nonzero obstruction
to a subsequent relative degree-three repair.
\end{proposition}

\begin{proof}
We use the signs of (AL.57)--(AL.62).  The letters below
refer to generators in the resolved complex $W$, before
tensoring with the single normalized generator of
$K(2)_2$.  Their unshifted weights are indicated by
their resolution subscripts.  Put
\[
\begin{gathered}
 C=[v_{g_3}],\qquad H=[u_{h_4}],\qquad K=[v_{g_4}],
       \qquad TC=2(K-H),\\
 x=v_{r_2},\quad y=v_{r_3},\quad
 g=u_{g_3},\quad l=u_{g_4},\\
 a=v_{e_1},\quad b=v_{e_2},\quad
 r=u_{r_2},\quad \sigma=u_{r_3},\quad
 E=u_{e_1},\quad F=u_{e_2}.
\end{gathered}                                      \tag{S09.98}
\]
Here $C,H,K$ are zero-forms; $x,y,g,l$ are one-forms;
$a,b,r,\sigma$ are two-forms; and $E,F$ are three-forms.
The weights of $C,H,K$ are $s+3,s+4,s+4$.
With $U=T^{[2]}$, the untwisted differentials are
\[
\begin{aligned}
 dC&=2g+Tx-2y,& dK&=2l+Ux-Ty,& dH&=-Tg+2l,\\
 dx&=2r+Ta-2b,&dy&=2\sigma+Ua-Tb,\\
 dg&=-Tr+2\sigma,&dl&=-Ur+T\sigma,\\
 da&=2E,&db&=2F,&dr&=-TE+2F,&
 d\sigma&=-UE+TF,\\
 dE&=0,&dF&=0.
\end{aligned}                                      \tag{S09.99}
\]
Every displayed generator acquires the same shift $s$;
coefficients in $S$ retain their original weights.
These formulas follow directly from the actual free
resolution and its nonnegative truncation.  In particular
the relation among $C,H,K$ is imposed in $W^0$ itself.

Define the integral forms
\[
 \begin{split}
 B&=aH-Cr-gx,\\
 \zeta&=EK-FC+ry-\sigma x-gb+la.
 \end{split}                                       \tag{S09.100}
\]
Both have weight $2s+5$; their form degrees are two and
three respectively.  Applying (S09.99), with the odd
sign in $d(gx)$, gives the exact identity
\[
                         dB=2\zeta.                \tag{S09.101}
\]
For example the zero-form/three-form part is
$2EH+TCE=2EK$, by (S09.98).  The terms $Tag$,
$Trx$, and $2gr$ cancel in pairs.  Torsion-freeness
then gives $d\zeta=0$.

We prove that $\zeta$ is not a boundary using the full
integral lattice in this arity and weight.  Let
\[
 E_W=\operatorname{im}\left(
 \operatorname{FreePD}_{S}(W)\longrightarrow
 \operatorname{FreePD}_{S_{\mathbb Q}}(W_{\mathbb Q})
 \right)
\]
with the zero twist, and give every $W$ generator arity
one.  Its arity-two, form-two group in weight $2s+5$
has the following $R$-basis:
\[
                  Cr,\quad gx,\quad aH,\quad aK,
                         \quad bC,\quad xy.         \tag{S09.102}
\]
Indeed the zero-form weights start with $s+3,s+4$,
with bases $C$ and $H,K$.  The one-form weights start
with $s+2$ on $x$, and $s+3$ on $g,y$.
The only additional apparent product $TaC$ equals
$2a(K-H)$.  The odd square $x^2$ is zero.  The
displayed six elements are rationally independent in
the free rational algebra, which proves the asserted
integral basis after taking the torsion-free image.

In the arity-two, form-three group, the summand with
one zero-form and one three-form has $R$-basis
\[
                            EH,\quad EK,\quad FC.
 \tag{S09.103}
\]
It is independent from the summand consisting of a
one-form times a two-form: the latter has a different
pair of form degrees in its rational free monomials.
Consequently the coefficient of $EK$ in (S09.103),
reduced modulo two, defines an actual $R$-linear
functional on the entire form-three lattice.
On the differentials of the ordered basis (S09.102)
its unreduced values are
\[
                         (-2,0,0,2,0,0),            \tag{S09.104}
\]
whereas its value on $\zeta$ is one.  This proves
nonboundary without any rational rank inference.
Together with (S09.101), it proves exact order two.

We now justify persistence through every stage of
saturation, including its later PD closures.  There
is an entirely algebraic version of the same operation
on $E_W$, using its augmentation
$\epsilon:E_W\to S$ in place of a vertex: inside
$(E_W)_{\mathbb Q}$, adjoin every homogeneous zero-form
$f$ with $df$ in the current integral algebra and
$\epsilon(f)\in S$, take PD closure, and iterate.
All stages preserve arity.  This follows by decomposing
$f$ into its finitely many arity components, since
the differential preserves arity, and then using the
addition and homogeneity identities for divided powers.

No new arity-one zero-form is introduced at any stage.
Its derivative lies in $W^1$, and the integral
exactness $H^1(W)=H^0(W)=0$ implies that its unique
rational primitive already lies in $W^0$.
All rational zero-forms of arity at least two have
weight at least $2(s+3)=2s+6$.  The arity-zero algebra
is the coefficient ring, and the augmentation condition
does not enlarge it.  It follows by induction on the
stages that the algebraic saturation agrees with $E_W$
in every weight at most $2s+5$, in every form degree.
This argument also allows the operation to be restricted
to any prescribed sufficiently large weight range.

In simplex two, the actual zero-twist algebra over
$G_R$ maps to $E_W$ by evaluating $G_R$ at vertex zero
and retaining the new $W$ generators.  Conversely the
coefficient inclusion gives a strict PD cochain map
$E_W$ into that simplex algebra, with composite the
identity.  For the persistence assertion only the first
of these maps is needed.  Its composition with
$\epsilon$ is exactly vertex-zero restriction of the
actual simplicial algebra, since $K(2)_0=0$.
Thus it sends each based saturation generator to a
generator allowed by the algebraic saturation, and
it sends the entire subsequent PD closure there as
well.  Induction gives an actual cochain map from the
saturated simplex algebra to the algebraic saturation
of $E_W$.  The image of $\zeta$ is the nonboundary
just proved.  This proves its persistence.  No
simplicial vertex-evaluation map $G_R\to S$ is claimed
or needed.

For the independent quadratic twist, use the old
identity disk described in the statement.  As in A-L11,
the normalized expressions $\gamma_2(y_0)$ and
$x_0y_0$ define actual linear maps from the sphere
parameter $K(2)$, and coefficient multiplication
defines the required maps on $L\otimes K(2)$.
In simplex two, the only changes in (S09.99) relevant
to (S09.100) are
\[
 da=2E-T\eta_0,\quad db=2F-U\eta_0,
 \qquad dE=Tz_0,\quad dF=Uz_0.                     \tag{S09.105}
\]
There is an integral zero-form $J\in W^0$ of weight
$s+5$ with
\[
              TH=2J,\qquad dJ=-Ug+Tl.
 \tag{S09.106}
\]
To see this without choosing any omitted resolution
generator, calculate $d(-Ug+Tl)=0$ using (S09.99) and
$T^2=2U$.  The definition
$W^0=\ker(W^1\to W^2)$ supplies the unique $J$ with
that derivative.  Since $d:W^0\to W^1$ is injective,
the equality $TdH=2(-Ug+Tl)$ proves $TH=2J$.
The correction is therefore an actual integral element,
not a rational generator added by assumption.

Put
\[
                         \zeta_{\mathrm{tw}}
                                  =\zeta-J\eta_0.
 \tag{S09.107}
\]
Equations (S09.101) and (S09.105)--(S09.106) now give
$dB=2\zeta_{\mathrm{tw}}$.  Hence this corrected
form is closed in the torsion-free algebra.  The strict
map that kills the old identity disk sends it to the
zero-twist $\zeta$.  Composing, in simplex two, with
the coefficient evaluation used above gives the same
map to $E_W$, compatible with the augmentation and
with every based saturation stage.  The corrected
class is therefore nonzero after saturation and has
exact order two there.

The actual resolution-augmentation comparison of AT16
sends the generators
\[
                 C,H,K,x,y,g,l,r,\sigma,J
\]
to zero.
It consequently sends $B$ and $\zeta_{\mathrm{tw}}$
to zero.  Its target is the two-term algebra of
A-S09-21, which has coefficient $H^0$ and zero $H^1$
in every weight, so A-SAT01 extends this very comparison
through the entire saturation.  In simplex two its
parameter has exactly the single independent generator,
so its target has $H^3=0$ there by A-S09-21.  This proves the stated failure
of the actual saturated comparison.
\end{proof}

The finite certificate
\texttt{code/A\_cosimplicial\_saturated\_H3\_certificate.py}
checks all thirty generator $d^2$ identities in the
untwisted and twisted algebras, both identities for
$dB$, both closedness identities, (S09.106), and the
complete six-column detector (S09.104).  Its exact
output is
\texttt{logs/A\_cosimplicial\_saturated\_H3\_certificate.json}.
The run took $0.014463970001088455$ seconds.
It also records the low normalized quadratic parameter
matrices used in the following calculation; persistence
through infinitely many saturation stages is proved
above and is not inferred from this finite computation.

\paragraph{The parameter calculation retained for the next repair.}
Let $e$ generate $K(\mathbb Z,2)_2$ and put
$a=s_2e,b=s_1e,c=s_0e$ in simplicial degree three.
After quotienting by degenerate terms, the degree-three
basis for both $\operatorname{Sym}^2 K(2)$ and
$\Gamma^2 K(2)$ is $ab,ac,bc$.  The boundary matrices
to their degree-two generators are respectively
$(1,0,-1)$ and $(2,0,-2)$.  Thus
$H_2\operatorname{Sym}^2 K(2)=0$ and
$H_2\Gamma^2 K(2)=\mathbb Z/2$.
The actual diagonal map
$\Gamma^2 K(2)\to K(2)\otimes K(2)$ sends
$\gamma_2(e)$ to $e\otimes e$, and
$\partial(b\otimes a)=e\otimes e$.
In particular it kills this degree-two homology class.
The expression $B(e)$ itself factors through that
diagonal by the bilinear formula
$B(t,u)=a(t)H(u)-C(t)r(u)-g(t)x(u)$.
No division of its off-diagonal differential by two
has been supplied, so this factorization alone does
not establish a relative period calculation for the
whole pair $(B,\zeta_{\mathrm{tw}})$.
